\documentclass{article}
\usepackage{graphicx} % Required for inserting images
\usepackage[spanish]{babel}
\usepackage{blindtext}
\usepackage{hyperref}


\title{Ryan Condal el \textit{"masiosare"} de nuestros tiempos}
\author{Adrián López Méndez}
\date{\today}


\begin{document}

\maketitle

\begin{figure}[h]
        \centering
        \includegraphics[width=1\linewidth]{dd.png}
        \label{fig:placeholder}
    \end{figure}
    
\section{El incidente del 2019}
Despues el rotundo fracaso que fue la temporada final de Juego de Tronos, historia procucida por Hbo en el 2011, adaptacion de la saga de libros \textit{''A song of ice and fire"} de George R.R Martin, el cual en gran parte fue por la falta de material que adaptar pues nuestro querido Jorge no habia terminado los libros para ese momento, quedandose en \textit{Vientos de invierno} y a dia de hoy seguimos en \textit{Vientos de invierno}.

Teniamos ya el preambulo de lo que puede desencadenar la falta de comunicacion entre los Showrunners y el señor Jorge, aunado a eso la exigencia de hbo por terminar su novena sinfonia, eso desenboco en la catastrofe que fue coronar a \textit{Bran ''EL roto''}; pero bueno 2019 paso y con eso llegaron buenas noticias a los fans, adaptarian otro libro interesante del mundo de hielo y fuego para Hbo proyecto con el nombre \textbf{''House Of The Dragon''}

\begin{figure}[h]
        \centering
        \includegraphics[width=0.5\linewidth]{image.png}
        \caption{Arciano}
        \label{fig:placeholder}
    \end{figure}

\section{¿Qué mató a HOTD?}
La catastrofe termino en mayo del 2019... \vspace{\baselineskip}

Y en octubre del mismo año se anunciaba el nuevo proyecto y la fecha de estreno oficial llegaria 3 años despues en el 2022, los fans estabamos emocionados por la noticia pero a la vez con miedo pues como dijo una persona muy sabia \textit{''Perdono, pero no olvido''} \textbf{-Gomita 2024}, aun asi eramos mas los ilusionados por un nuevo proyecto de las manos de George y de el gran Judas, el showrunner Ryan Condal, el cual expresó su profundo respeto y compromiso hacia el autor George R. R. Martin, dejando claro que era fanático de Canción de Hielo y Fuego desde hacía casi 25 años y que trabajar con él era el mayor privilegio de su vida.

\begin{figure}[h]
        \centering
        \includegraphics[width=0.6\linewidth]{i-would-prefer-that-our-dysfunction-stays-behind-closed-door_gg49.png}
        \caption{George R.R}
        \label{fig:placeholder}
    \end{figure}
    
En sus declaraciones iniciales y durante la producción de la primera temporada, Condal se comprometió firmemente a incluir al autor en el proceso creativo, asegurando que mantendrían una estrecha colaboración en la adaptación de los borradores y guiones, y asi fue, excelente y consolador durante la primera temporada de la serie, tan bello que volviamos a sentirnos fieles a la madre de dragones, rompedora de cadenas, fieles almirantes del bando Negro despues de la gatada que hizo Alicent Hightower al reino, pero ese sentimiento duro poco, pues los problemas no se hiceron esperar.

\section{La historia se repite}

 Apesar de aquellas primeras promesas de fidelidad al material original y respeto al autor, la relación sufrió un gran quiebre tras el estreno de la segunda temporada. En entrevistas y declaraciones ambos expusieron perspectivas opuestas sobre el rumbo de la serie.
 
 \begin{figure}[h]
        \centering
        \includegraphics[width=1\linewidth]{matt-smith.png}
        \caption{\textbf{POV}: El fandom leyedo las declaraciones de Jorjais}
        \label{fig:placeholder}
    \end{figure}
    
George R. R. Martin reveló públicamente que la relación se volvió "pésima". Manifestó que Condal "básicamente dejó de escucharlo" en la segunda temporada e ignoró sus advertencias sobre los cambios en la trama de su libro Fuego y sangre, culminando en la célebre y amarga frase del escritor: \textit{''Esta ya no es mi historia''}.\vspace{\baselineskip}

A lo que Ryan desacredito a Jorjito argumentando que el le sabe a los libros pero no a la tele, y es asi como una serie con tanto potencial termino cometiendo los mismos errores que su predecesora, estando destinada al mismo fracaso.\vspace{\baselineskip}

Aun con todo eso, el mundo de Hielo y Fuego es tan basto y fantastico que la serie no termino siendo una basura como su hermana mayor, logro mantenerse y entregar capitulos con guiones aceptables, actuaciones estelares y un hilo cada vez mas enredado que solo Ryan sabra como terminar su despapaye.\vspace{\baselineskip}

Hubieron tantas historias que no puedieron tener el impacto que debieron por lo mal escritas que estuvieron para la television, como el relato de sangre y queso para la muerte del hijo de Aegon ''El usurpador"  o la relacion inventada de Mysaria y Rhaenyra, las muertes "accidentales" de los hijos de la reina negra, La batalla del Gaznate y toda la trama de Harrenhal son tan mediocres que no terminaria nunca de expresar mi odio por Rhyan Condal. \vspace{\baselineskip}

Nuestras esperanzas aun recaen en que Don George saque pronto vientos de inviernos sin mas... \vspace{\baselineskip}

\textit{VALLAR DOHAERIS}

\section{Referencias}
\begin{itemize}
  \item \href{https://www.bbc.com/news/entertainment-arts-50229837}{BBC News. (2019, 30 octubre). Game of Thrones prequel House of the Dragon ordered by HBO. https://www.bbc.com/news/entertainment-arts-50229837}
  \item \href{https://www.imdb.com/es/title/tt0944947/episodes/?season=8}{Game of Thrones (Serie de TV 2011–2019) - Lista de episodios - IMDb. (n.d.). IMDb. https://www.imdb.com/es/title/tt0944947/episodes/?season=8}
  \item \href{https://mashable.com/article/george-r-r-martin-house-of-the-dragon-season-3-ryan-condal-abysmal}{Edwards, B. (2026, January 16). George R.R. Martin talks “abysmal” relationship with “House of the Dragon” showrunner Ryan Condal. Mashable. https://mashable.com/article/george-r-r-martin-house-of-the-dragon-season-3-ryan-condal-abysmal}
  \item \href{https://es.overleaf.com/learn/latex/Learn_LaTeX_in_30_minutes}{Learn LaTeX in 30 minutes. (s.f.). Overleaf, Editor de LaTeX Online. https://es.overleaf.com/learn/latex/Learn\_LaTeX\_in\_30\_minutes}
\end{itemize}

\end{document}
